% reality/km_line_and_stationing-DE.tex
\chapter{Kilometrierungslinie und Stationierung}
\label{chap:km-line-and-stationing}
\index{Bezugssystem!Stationierung}

Stationierung ist die operative Adressierungsebene von Eisenbahnsystemen. Sie
ist weder mit konstruktiver Alignment-Identität noch mit metrischer Einbettung
oder realisierter Geometrie identisch.

\section{Rolle der Stationierung}
\begin{definition}[Kilometrierungslinie]
\index{Bezugssystem!Kilometrierungslinie}
Eine Kilometrierungslinie ist eine longitudinale Referenzstruktur, die in
einem Routenkontext Stationierungswerte zuordnet.
\end{definition}
\begin{definition}[Stationierungsebene]
\index{Bezugssystem!Stationierungsebene}
Die Stationierungsebene bildet Eisenbahnpositionen auf operative Adresswerte
ab.
\end{definition}

Stationierung beantwortet, wie eine Infrastrukturposition adressiert wird;
Geometrie beantwortet, welche Form realisiert ist; CRS beantworten, wo die
Einbettung ausgedrückt wird.

\section{Bezug zur Geometrie}
Eine Kilometrierungslinie kann mit einer physischen Gleismittellinie
zusammenfallen, muss es aber nicht. In mehrgleisigen und
Verzweigungskontexten können mehrere physische Alignments auf dieselbe
Stationierungsstruktur Bezug nehmen.

\begin{definition}[Äußere Stationierung]
\index{Bezugssystem!äußere Stationierung}
Äußere Stationierung beschreibt die Abbildung zwischen einer physischen
Gleis-Alignment-Koordinate und einer externen Kilometrierungsreferenz.
\end{definition}
Sie ist erforderlich, wenn Stationierungsreferenz und physische Gleisreferenz
nicht zusammenfallen.

\section{Trennung der Ebenen}
\begin{proposition}
Eisenbahnmodellierung erfordert die Trennung
\[
\text{räumliche Einbettung}\neq\text{geometrisches Alignment}
\neq\text{Stationierungsreferenz}.
\]
\end{proposition}
Die Vermengung dieser Ebenen erzeugt systematische Fehler in Interpretation,
Netzadressierung und Ziel--Realisiert-Vergleich.

\section{Topologie und operative Kontinuität}
\index{Topologie!operative Kontinuität}
Kilometrierungsreferenzen stellen operative Kontinuität über Gleiswechsel,
Verzweigungen, Bahnhofsbereiche und Routenübergänge hinweg bereit, selbst wenn
geometrische Kontinuität durch mehrere physische Alignments dargestellt wird.

\section{Verbindung zu AIM}
\begin{summary}
In AIM ist Stationierung eine ausdrückliche Referenzebene für operative
Adressierung und Vergleichsabläufe. Sie bleibt von konstruktiver Identität,
CRS-Einbettung und realisierter Gleisgeometrie verschieden.
\end{summary}
